\documentclass[11pt, a4paper]{article}
\usepackage[a4paper, margin=1.8cm]{geometry}
\usepackage[T2A]{fontenc}
\usepackage[utf8]{inputenc}
\usepackage[russian,english]{babel}
\usepackage{enumitem}
\setlist[itemize]{label=-}
\usepackage{amsmath}
\usepackage{booktabs}
\usepackage{array}
\usepackage{hyperref}
\usepackage{caption}
\usepackage{adjustbox}

\setcounter{secnumdepth}{0}

% Подпись таблицы: справа "Таблица N", ниже по центру — описание
\captionsetup[table]{
  labelsep=newline,
  justification=raggedright,
  singlelinecheck=false,
  labelfont=normalfont,
  textfont=normalfont,
  labelformat=simple,
  labelsep=space
}
\renewcommand{\thetable}{\arabic{table}}
\DeclareCaptionLabelFormat{russian}{Таблица~#2}
\captionsetup[table]{labelformat=russian}

% Хак: метка справа, описание по центру под ней
\makeatletter
\long\def\@makecaption#1#2{%
  \vskip\abovecaptionskip
  \begin{flushright}\normalfont #1\end{flushright}%
  \begin{center}\normalfont #2\end{center}%
  \vskip\belowcaptionskip}
\makeatother

\begin{document}

\tableofcontents
\newpage

\section{Цель}
Микропрограммная реализация и исследование учебной ЭВМ с заданной архитектурой.

\section{Задание}
Определить архитектуру ЭВМ, система команд которой состоит из одноадресных команд, использующих прямую адресацию; разработать структурную схему и алгоритм работы ЭВМ; составить и отладить микропрограмму командного цикла ЭВМ; составить и выполнить программу операции сдвига 16-разрядного двоичного кода $X$ на $N$ разрядов в сторону старших разрядов (сдвиг логический).

\section{Определение архитектуры и программирование}

\subsection{Форматы данных}
Обрабатываемые данные и команды представляются 16-разрядным двоичным кодом. В машинной команде старшие 8 разрядов, разряды 15--8, определяют код операции, а младшие 8 разрядов, разряды 7--0, задают прямой адрес ячейки оперативной памяти $A$.

\subsection{Программно-доступные регистры}
\begin{itemize}
	\item \textbf{AC} --- аккумулятор;
	\item \textbf{PC} --- программный счётчик команд;
	\item \textbf{PZ} --- программный признак нуля.
\end{itemize}

\subsection{Система команд}
Система команд учебной ЭВМ представлена в таблице 1, где $M[A]$ --- ячейка оперативной памяти с адресом $A$. Знак «+» в столбце признака PZ означает, что по результату выполнения данной команды устанавливается новое значение признака нуля, а знак «-» свидетельствует о сохранении прежнего значения признака.

\begin{table}[h!]
	\centering
	\caption{Система команд ЭВМ}
	\begin{tabular}{|l|c|l|c|}
		\hline
		\textbf{Наименование} & \textbf{Мнемоника} & \textbf{Описание}                                 & \textbf{PZ} \\
		\hline
		Очистка               & \texttt{CLM A}     & $M[A] := 00..0$, $PC := PC + 1$                   & -           \\
		\hline
		Загрузка AC           & \texttt{LDA A}     & $AC := M[A]$, $PC := PC + 1$                      & +           \\
		\hline
		Запись AC             & \texttt{MOV A}     & $M[A] := AC$, $PC := PC + 1$                      & -           \\
		\hline
		Декремент             & \texttt{DEC A}     & $M[A] := M[A] - 1$, $PC := PC + 1$                & +           \\
		\hline
		Логическое И          & \texttt{AND A}     & $AC := AC \wedge M[A]$, $PC := PC + 1$            & -           \\
		\hline
		Логическое ИЛИ        & \texttt{OR A}      & $AC := AC \vee M[A]$, $PC := PC + 1$              & -           \\
		\hline
		Сдвиг влево           & \texttt{SHL A}     & $M[A] := 2 \cdot M[A]$, $PC := PC + 1$            & -           \\
		\hline
		Безусловный переход   & \texttt{JMP A}     & $PC := A$                                         & -           \\
		\hline
		Переход, если ноль    & \texttt{JZ A}      & Если $PZ = 1$, то $PC := A$, иначе $PC := PC + 1$ & -           \\
		\hline
		Нет операции          & \texttt{NOP}       & $PC := PC + 1$                                    & -           \\
		\hline
		Останов               & \texttt{HLT A}     & $PC := A$, останов                                & -           \\
		\hline
	\end{tabular}
\end{table}
\newpage
\subsection{Программа}
Программа выполнения логического сдвига 16-разрядного числа $X$ на $N$ разрядов влево на языке мнемоник приведена ниже.

\begin{center}
	\begin{tabular}{|l|l|l|}
		\hline
		\textbf{Метка} & \textbf{Команда} & \textbf{Назначение команды}                          \\
		\hline
		               & \texttt{LDA AX}  & Загрузка исходного числа $X$ в аккумулятор AC        \\
		               & \texttt{MOV AR}  & Запись значения из AC в ячейку результата AR         \\
		               & \texttt{LDA N1}  & Загрузка количества сдвигов $N$ в AC                 \\
		               & \texttt{MOV AN1} & Запись значения $N$ в рабочий счётчик цикла AN1      \\
		\texttt{M1:}   & \texttt{NOP}     & Метка начала цикла сдвига                            \\
		               & \texttt{JZ M2}   & Если $N1 = 0$, то переход на выход к метке M2        \\
		               & \texttt{SHL AR}  & Логический сдвиг числа в ячейке AR на 1 разряд влево \\
		               & \texttt{DEC AN1} & Вычитание единицы из счётчика AN1                    \\
		               & \texttt{JMP M1}  & Безусловный переход на метку M1                      \\
		\texttt{M2:}   & \texttt{HLT SA}  & Загрузка стартового адреса SA в PC и останов         \\
		\hline
	\end{tabular}
\end{center}

В программе приняты следующие символические обозначения адресов ячеек памяти:
\begin{itemize}
	\item \textbf{SA} --- ячейка памяти с адресом 00, хранящая начальный адрес запуска программы;
	\item \textbf{AX} --- адрес ячейки 01, в которой находится исходное сдвигаемое число $X$;
	\item \textbf{N1} --- адрес ячейки 02, в которой находится число $N$, определяющее количество позиций сдвига;
	\item \textbf{AR} --- адрес ячейки 03, в которой формируется и хранится результат логического сдвига $R$;
	\item \textbf{AN1} --- адрес ячейки 04, используемой в качестве изменяемого рабочего счётчика цикла $N1$.
\end{itemize}

\section{Кодирование программы и распределение памяти программ и данных}
Наименования команд, мнемонические обозначения и коды операций в двоичной и шестнадцатеричной системах счисления приведены в таблице 2.

\begin{table}[h!]
	\centering
	\caption{Коды операций учебной ЭВМ}
	\begin{tabular}{|l|c|c|c|}
		\hline
		\textbf{Наименование} & \textbf{Мнемоника} & \textbf{Код операции (2сс)} & \textbf{Код операции (16сс)} \\
		\hline
		Останов               & \texttt{HLT A}     & \texttt{00000000}           & \texttt{00}                  \\
		Очистка               & \texttt{CLM A}     & \texttt{00000001}           & \texttt{01}                  \\
		Загрузка AC           & \texttt{LDA A}     & \texttt{00000010}           & \texttt{02}                  \\
		Запись AC             & \texttt{MOV A}     & \texttt{00000011}           & \texttt{03}                  \\
		Декремент             & \texttt{DEC A}     & \texttt{00000100}           & \texttt{04}                  \\
		Логическое И          & \texttt{AND A}     & \texttt{00000101}           & \texttt{05}                  \\
		Логическое ИЛИ        & \texttt{OR A}      & \texttt{00000110}           & \texttt{06}                  \\
		Сдвиг влево           & \texttt{SHL A}     & \texttt{00000111}           & \texttt{07}                  \\
		Безусловный переход   & \texttt{JMP A}     & \texttt{00001000}           & \texttt{08}                  \\
		Переход, если ноль    & \texttt{JZ A}      & \texttt{00001001}           & \texttt{09}                  \\
		Нет операции          & \texttt{NOP}       & \texttt{00001010}           & \texttt{0A}                  \\
		\hline
	\end{tabular}
\end{table}

Распределение оперативной памяти для хранения исходных данных, ячейки 00--04, и машинных команд программы сдвига, ячейки 05--0E, приведено в таблице 3. В качестве контрольного примера заданы исходное число $X = 8006_{16}$ и число сдвигов $N = 0003_{16}$; после выполнения в ячейке результата $R$ формируется значение $0030_{16}$.

\begin{table}[h!]
	\centering
	\caption{Распределение памяти программ и данных}
	\begin{tabular}{|c|c|c|l|l|}
		\hline
		\textbf{Адрес} & \textbf{Код (2сс)}        & \textbf{Код (16сс)} & \textbf{Мнемоника}  & \textbf{Комментарий}                     \\
		\hline
		\texttt{00}    & \texttt{0000000000000101} & \texttt{0005}       & \texttt{SA}         & Начальный адрес программы 05             \\
		\texttt{01}    & \texttt{1000000000000110} & \texttt{8006}       & \texttt{X}          & Исходное сдвигаемое число                \\
		\texttt{02}    & \texttt{0000000000000011} & \texttt{0003}       & \texttt{N}          & Количество разрядов для сдвига           \\
		\texttt{03}    & \texttt{0000000000110000} & \texttt{0030}       & \texttt{RESULT}     & Результат сдвига, до пуска \texttt{0000} \\
		\texttt{04}    & \texttt{0000000000000000} & \texttt{0000}       & \texttt{N1}         & Изменяемый рабочий счётчик $N1$          \\
		\texttt{05}    & \texttt{0000001000000001} & \texttt{0201}       & \texttt{LDA X}      & Загрузка исходного числа $X$ в AC        \\
		\texttt{06}    & \texttt{0000001100000011} & \texttt{0303}       & \texttt{MOV R}      & Сохранение числа $X$ в ячейку $R$        \\
		\texttt{07}    & \texttt{0000001000000010} & \texttt{0202}       & \texttt{LDA N}      & Загрузка количества сдвигов $N$ в AC     \\
		\texttt{08}    & \texttt{0000001100000100} & \texttt{0304}       & \texttt{MOV N1}     & Сохранение $N$ в рабочий счётчик $N1$    \\
		\texttt{09}    & \texttt{0000101000000000} & \texttt{0A00}       & \texttt{M1: NOP}    & Метка начала цикла                       \\
		\texttt{0A}    & \texttt{0000100100001110} & \texttt{090E}       & \texttt{JZ M2}      & Если $N1 = 0$, переход на M2             \\
		\texttt{0B}    & \texttt{0000011100000011} & \texttt{0703}       & \texttt{SHL R}      & Логический сдвиг $R$ влево на 1 разряд   \\
		\texttt{0C}    & \texttt{0000010000000100} & \texttt{0404}       & \texttt{DEC N1}     & Декремент счётчика $N1$                  \\
		\texttt{0D}    & \texttt{0000100000001001} & \texttt{0809}       & \texttt{JMP M1}     & Безусловный переход на метку M1          \\
		\texttt{0E}    & \texttt{0000000000000101} & \texttt{0005}       & \texttt{M2: HLT SA} & Загрузка SA в PC и останов программы     \\
		\hline
	\end{tabular}
\end{table}

\section{Разработка структуры и алгоритма работы}
Алгоритм функционирования процессора строится на основе командного цикла, включающего четыре основные фазы.
\begin{itemize}
	\item \textbf{Начальный пуск, МК 00--03.} При запуске устройства в регистре константы RE путём сдвига вправо в цикле формируется 8-разрядная маска младшего байта $00FF_{16}$, МК 00--01. Затем из нулевой ячейки ОЗУ с адресом 00 считывается стартовый адрес программы SA и помещается в программный счётчик команд PC, регистр R6, МК 02--03.
	\item \textbf{Выборка и дешифрация команды, МК 04--05.} В МК 04 содержимое счётчика команд PC выдаётся в регистр адреса ОЗУ, $RA := PC$, а сам счётчик команд увеличивается на единицу, $PC := PC + 1$. В МК 05 считанная из ОЗУ команда логически умножается на маску RE, благодаря чему в регистр команд RF записывается чистый адрес операнда $A$, $RF := M[RA] \wedge 00FF_{16}$. Одновременно старший байт команды, код операции, с шины данных подаётся на вход памяти начальных адресов ПНА по микроинструкции JMAP, осуществляя переход к подмикропрограмме соответствующей операции.
	\item \textbf{Выполнение операции, МК 0B--28.} Подмикропрограмма выбранной машинной команды передаёт адрес операнда из RF в регистр адреса памяти, $RA := RF$, и выполняет требуемую обработку данных в ОЗУ и аккумуляторе R0. Команды, не изменяющие признак нуля PZ, а именно \texttt{CLM}, \texttt{MOV}, \texttt{AND}, \texttt{OR}, \texttt{SHL}, \texttt{JMP}, \texttt{JZ} и \texttt{NOP}, по завершении выполняют безусловный переход обратно на МК 04 к выборке следующей команды.
	\item \textbf{Формирование признака нуля PZ, МК 07--09.} Команды \texttt{LDA}, МК 0F, и \texttt{DEC}, МК 16, после выполнения операции переходят к блоку установки программного признака PZ. В МК 07 регистр признака R7 предварительно обнуляется, $R7 := 0$. В МК 08 значение аккумулятора R0 проверяется через аппаратный микропрограммный признак $Z$ АЛУ: если результат не равен нулю, $\overline{Z} = 1$, происходит переход на МК 04; если результат равен нулю, $Z = 1$, в МК 09 значение R7 увеличивается на единицу, $R7 := 1$, после чего управление также возвращается на МК 04.
\end{itemize}

\section{Микропрограммная реализация ЭВМ}
Для реализации архитектуры учебной ЭВМ на базе микропрограммируемого секционного микропроцессора выполнено распределение регистров регистрового запоминающего устройства, представленное в таблице 4.

	\begin{table}[h!]
		\centering
		\caption{Распределение регистров РЗУ и операционного устройства}
		\begin{adjustbox}{max width=\textwidth}

		\begin{tabular}{|c|c|l|}
			\hline
			\textbf{Регистр РЗУ} & \textbf{Обозначение в ЭВМ} & \textbf{Назначение регистра в микропрограмме}                                 \\
			\hline
			\texttt{R0}          & \texttt{AC}                & Аккумулятор                                                                   \\
			\texttt{R1 -- R5}    & ---                        & Свободные регистры РЗУ                                                        \\
			\texttt{R6}          & \texttt{PC}                & Программный счётчик команд, адрес текущей команды в ОЗУ                       \\
			\texttt{R7}          & \texttt{PZ}                & Регистр программного признака нуля, \texttt{0001} ноль, \texttt{0000} не ноль \\
			\texttt{R8 -- RD}    & ---                        & Свободные регистры РЗУ                                                        \\
			\texttt{RE}          & Константа $00FF_{16}$      & Регистр маски для выделения младшего байта, адреса операнда $A$               \\
			\texttt{RF}          & \texttt{RK[A]}             & Регистр адреса операнда команды, адрес ячейки памяти $A$                      \\
			\texttt{RA}          & Адрес ЗУ                   & Регистр адреса оперативного запоминающего устройства                          \\
			\hline
		\end{tabular}
	\end{adjustbox}
	\end{table}

Соответствие между кодами операций машинных команд, индексами ячеек таблицы ПНА, и начальными адресами соответствующих подмикропрограмм в управляющей памяти приведено в таблице 5.

\begin{table}[h!]
	\centering
	\caption{Коды операций и начальные адреса подмикропрограмм, таблица ПНА}
	\begin{adjustbox}{max width=\textwidth}
	\begin{tabular}{|l|c|c|c|c|}
		\hline
		\textbf{Мнемоника} & \textbf{Код операции (2сс)} & \textbf{Адрес ПНА (16сс)} & \textbf{Адрес первой МК (2сс)} & \textbf{Адрес первой МК (16сс)} \\
		\hline
		\texttt{HLT A}     & \texttt{00000000}           & \texttt{00}               & \texttt{00101000}              & \texttt{28}                     \\
		\texttt{CLM A}     & \texttt{00000001}           & \texttt{01}               & \texttt{00001011}              & \texttt{0B}                     \\
		\texttt{LDA A}     & \texttt{00000010}           & \texttt{02}               & \texttt{00001110}              & \texttt{0E}                     \\
		\texttt{MOV A}     & \texttt{00000011}           & \texttt{03}               & \texttt{00010001}              & \texttt{11}                     \\
		\texttt{DEC A}     & \texttt{00000100}           & \texttt{04}               & \texttt{00010100}              & \texttt{14}                     \\
		\texttt{AND A}     & \texttt{00000101}           & \texttt{05}               & \texttt{00011000}              & \texttt{18}                     \\
		\texttt{OR A}      & \texttt{00000110}           & \texttt{06}               & \texttt{00011011}              & \texttt{1B}                     \\
		\texttt{SHL A}     & \texttt{00000111}           & \texttt{07}               & \texttt{00011110}              & \texttt{1E}                     \\
		\texttt{JMP A}     & \texttt{00001000}           & \texttt{08}               & \texttt{00100100}              & \texttt{24}                     \\
		\texttt{JZ A}      & \texttt{00001001}           & \texttt{09}               & \texttt{00100011}              & \texttt{23}                     \\
		\texttt{NOP}       & \texttt{00001010}           & \texttt{0A}               & \texttt{00100110}              & \texttt{26}                     \\
		\hline
	\end{tabular}
	\end{adjustbox}
\end{table}

\section{Кодирование микропрограммы командного цикла}
В таблице 6 приведено полное кодирование микропрограммы командного цикла учебной ЭВМ во вкладке «ОУ и УУ».

	{\scriptsize
		\begin{table}[h!]
			\centering
			\caption{Микропрограмма командного цикла учебной ЭВМ}
				\begin{adjustbox}{max width=\textwidth}

			\begin{tabular}{|c|c|c|c|c|c|c|c|c|c|c|c|c|c|c|c|c|c|l|}
				\hline
				\textbf{№}  & \textbf{МИ}                       & \multicolumn{2}{c|}{\textbf{РЗУ}} & \multicolumn{3}{c|}{\textbf{Упр. АЛУ}} & \multicolumn{3}{c|}{\textbf{Упр. ОЗУ}} & \textbf{Шина} & \textbf{МИ} & \multicolumn{3}{c|}{\textbf{Упр. усл.}} & \multicolumn{3}{c|}{\textbf{Упр. УУ}} & \textbf{Выполняемая}                                                                                                                                                       \\
				\textbf{МК} & \textbf{I8-0}                     & \textbf{A}                        & \textbf{B}                             & \textbf{C0}                            & \textbf{OE}   & \textbf{SC} & \textbf{CS}                             & \textbf{W}                            & \textbf{EA}          & \textbf{D11-0} & \textbf{I3-0} & \textbf{A} & \textbf{U} & \textbf{CCE} & \textbf{CI} & \textbf{RLD} & \textbf{OE} & \textbf{микрооперация}         \\
				\hline
				00          & 571                               & E                                 & E                                      & 0                                      & 0             & 00          & 1                                       & 1                                     & 1                    & 006            & C             & 00         & 0          & 0            & 1           & 1            & 0           & СЦ := 6, RE := 7FFF            \\
				01          & 533                               & 0                                 & E                                      & 0                                      & 0             & 00          & 1                                       & 1                                     & 1                    & 001            & 9             & 00         & 0          & 0            & 1           & 1            & 0           & Сдвиг RE, RE := 00FF           \\
				02          & 143                               & 0                                 & 6                                      & 0                                      & 0             & 00          & 1                                       & 1                                     & 0                    & 000            & E             & 00         & 0          & 0            & 1           & 1            & 0           & RA := 0, адрес SA              \\
				03          & 337                               & 0                                 & 6                                      & 0                                      & 1             & 00          & 0                                       & 1                                     & 1                    & 000            & E             & 00         & 0          & 0            & 1           & 1            & 0           & R6, PC := M[00]                \\
				04          & 203                               & 6                                 & 6                                      & 1                                      & 0             & 00          & 1                                       & 1                                     & 0                    & 000            & E             & 00         & 0          & 0            & 1           & 1            & 0           & RA := R6, R6 := R6 + 1         \\
				05          & 245                               & E                                 & F                                      & 0                                      & 1             & 00          & 0                                       & 1                                     & 1                    & 000            & 2             & 00         & 0          & 0            & 1           & 1            & 0           & RF := M[RA] $\wedge$ RE, JMAP  \\
				06          & \multicolumn{18}{l|}{Разделитель}                                                                                                                                                                                                                                                                                                                                                                                                                    \\
				07          & 343                               & 0                                 & 7                                      & 0                                      & 1             & 00          & 1                                       & 1                                     & 1                    & 000            & E             & 00         & 0          & 0            & 1           & 1            & 0           & R7 := 0, сброс PZ              \\
				08          & 133                               & 0                                 & 0                                      & 0                                      & 1             & 00          & 1                                       & 1                                     & 1                    & 004            & 3             & 00         & 0          & 0            & 1           & 1            & 0           & Если R0 $\neq$ 0, на МК 04     \\
				09          & 303                               & 0                                 & 7                                      & 1                                      & 1             & 00          & 1                                       & 1                                     & 1                    & 004            & 3             & 00         & 0          & 1            & 1           & 1            & 0           & R7 := R7 + 1, на МК 04         \\
				0A          & \multicolumn{18}{l|}{Разделитель}                                                                                                                                                                                                                                                                                                                                                                                                                    \\
				0B          & 133                               & 0                                 & F                                      & 0                                      & 0             & 00          & 1                                       & 1                                     & 0                    & 000            & E             & 00         & 0          & 0            & 1           & 1            & 0           & CLM A: RA := RF[A]             \\
				0C          & 143                               & 0                                 & F                                      & 0                                      & 0             & 00          & 0                                       & 0                                     & 1                    & 004            & 3             & 00         & 0          & 1            & 1           & 1            & 0           & M[RA] := 0, на МК 04           \\
				0D          & \multicolumn{18}{l|}{Разделитель}                                                                                                                                                                                                                                                                                                                                                                                                                    \\
				0E          & 133                               & 0                                 & F                                      & 0                                      & 0             & 00          & 1                                       & 1                                     & 0                    & 000            & E             & 00         & 0          & 0            & 1           & 1            & 0           & LDA A: RA := RF[A]             \\
				0F          & 337                               & 0                                 & 0                                      & 0                                      & 1             & 00          & 0                                       & 1                                     & 1                    & 007            & 3             & 00         & 0          & 1            & 1           & 1            & 0           & R0 := M[RA], на МК 07          \\
				10          & \multicolumn{18}{l|}{Разделитель}                                                                                                                                                                                                                                                                                                                                                                                                                    \\
				11          & 133                               & 0                                 & F                                      & 0                                      & 0             & 00          & 1                                       & 1                                     & 0                    & 000            & E             & 00         & 0          & 0            & 1           & 1            & 0           & MOV A: RA := RF[A]             \\
				12          & 134                               & 0                                 & 0                                      & 0                                      & 0             & 00          & 0                                       & 0                                     & 1                    & 004            & 3             & 00         & 0          & 1            & 1           & 1            & 0           & M[RA] := R0, на МК 04          \\
				13          & \multicolumn{18}{l|}{Разделитель}                                                                                                                                                                                                                                                                                                                                                                                                                    \\
				14          & 133                               & 0                                 & F                                      & 0                                      & 0             & 00          & 1                                       & 1                                     & 0                    & 000            & E             & 00         & 0          & 0            & 1           & 1            & 0           & DEC A: RA := RF[A]             \\
				15          & 337                               & 0                                 & 0                                      & 0                                      & 1             & 00          & 0                                       & 1                                     & 1                    & 000            & E             & 00         & 0          & 0            & 1           & 1            & 0           & R0 := M[RA]                    \\
				16          & 313                               & 0                                 & 0                                      & 0                                      & 0             & 00          & 0                                       & 0                                     & 1                    & 007            & 3             & 00         & 0          & 1            & 1           & 1            & 0           & M[RA] := R0 - 1, на 07         \\
				17          & \multicolumn{18}{l|}{Разделитель}                                                                                                                                                                                                                                                                                                                                                                                                                    \\
				18          & 133                               & 0                                 & F                                      & 0                                      & 0             & 00          & 1                                       & 1                                     & 0                    & 000            & E             & 00         & 0          & 0            & 1           & 1            & 0           & AND A: RA := RF[A]             \\
				19          & 345                               & 0                                 & 0                                      & 1                                      & 1             & 00          & 0                                       & 1                                     & 1                    & 004            & 3             & 00         & 0          & 1            & 1           & 1            & 0           & R0 := R0 $\wedge$ M[RA], на 04 \\
				1A          & \multicolumn{18}{l|}{Разделитель}                                                                                                                                                                                                                                                                                                                                                                                                                    \\
				1B          & 133                               & 0                                 & F                                      & 0                                      & 0             & 00          & 1                                       & 1                                     & 0                    & 000            & E             & 00         & 0          & 0            & 1           & 1            & 0           & OR A: RA := RF[A]              \\
				1C          & 335                               & 0                                 & 0                                      & 1                                      & 1             & 00          & 0                                       & 1                                     & 1                    & 004            & 3             & 00         & 0          & 1            & 1           & 1            & 0           & R0 := R0 $\vee$ M[RA], на 04   \\
				1D          & \multicolumn{18}{l|}{Разделитель}                                                                                                                                                                                                                                                                                                                                                                                                                    \\
				1E          & 133                               & 0                                 & F                                      & 0                                      & 0             & 00          & 1                                       & 1                                     & 0                    & 000            & E             & 00         & 0          & 0            & 1           & 1            & 0           & SHL A: RA := RF[A]             \\
				1F          & 337                               & 0                                 & 0                                      & 0                                      & 1             & 00          & 0                                       & 1                                     & 1                    & 000            & E             & 00         & 0          & 0            & 1           & 1            & 0           & R0 := M[RA]                    \\
				20          & 733                               & 0                                 & 0                                      & 0                                      & 1             & 00          & 1                                       & 1                                     & 1                    & 004            & E             & 00         & 0          & 0            & 1           & 1            & 0           & R0 := 2$\cdot$R0, сдвиг влево  \\
				21          & 133                               & 0                                 & 0                                      & 0                                      & 0             & 00          & 0                                       & 0                                     & 1                    & 004            & 3             & 00         & 0          & 1            & 1           & 1            & 0           & M[RA] := R0, на МК 04          \\
				22          & \multicolumn{18}{l|}{Разделитель}                                                                                                                                                                                                                                                                                                                                                                                                                    \\
				23          & 113                               & 0                                 & 7                                      & 0                                      & 1             & 00          & 1                                       & 1                                     & 1                    & 004            & 3             & 00         & 0          & 0            & 1           & 1            & 0           & JZ A: проверка R7 - 1          \\
				24          & 334                               & F                                 & 6                                      & 0                                      & 1             & 00          & 1                                       & 1                                     & 1                    & 004            & 3             & 00         & 0          & 1            & 1           & 1            & 0           & JMP A: R6 := RF, на 04         \\
				25          & \multicolumn{18}{l|}{Разделитель}                                                                                                                                                                                                                                                                                                                                                                                                                    \\
				26          & 133                               & 0                                 & 0                                      & 0                                      & 1             & 00          & 1                                       & 1                                     & 1                    & 004            & 3             & 00         & 0          & 1            & 1           & 1            & 0           & NOP: переход на МК 04          \\
				27          & \multicolumn{18}{l|}{Разделитель}                                                                                                                                                                                                                                                                                                                                                                                                                    \\
				28          & 334                               & F                                 & 6                                      & 0                                      & 1             & 00          & 1                                       & 1                                     & 1                    & 000            & 3             & 00         & 0          & 1            & 1           & 1            & 0           & HLT A: R6 := RF, на 00         \\
				\hline
			\end{tabular}
		\end{adjustbox}
		\end{table}
	}

\section{Экранные формы и результаты отладки}
В ходе отладки микропрограммы командного цикла и программы логического сдвига в имитаторе микропрограммируемого микропроцессора для контрольных исходных данных $X = 8006_{16}$ и $N = 0003_{16}$ получены следующие результаты.
\begin{itemize}
	\item В ячейке ОЗУ \texttt{03} сформировано сдвинутое значение $0030_{16}$, старшие разряды выдвинуты за пределы разрядной сетки, младшие заполнены нулями; при тестовом прогоне с $N = 0001_{16}$ результат равен $000C_{16}$.
	\item Рабочий счётчик цикла в ячейке \texttt{04} обнулён.
	\item Состояние регистров РЗУ после выполнения команды \texttt{HLT SA}: \texttt{R0} = \texttt{0000}, \texttt{R6} = \texttt{0005}, \texttt{R7} = \texttt{0001}, \texttt{RE} = \texttt{00FF}, \texttt{RF} = \texttt{0005}.
\end{itemize}

\section{Вывод}
В ходе выполнения лабораторной работы была определена архитектура учебной ЭВМ с одноадресными командами и прямой адресацией. Разработана система из 11 машинных команд, определены форматы данных и команд, выполнено распределение программно-доступных и рабочих регистров РЗУ: аккумулятор R0, программный счётчик R6, регистр программного признака нуля R7, регистр маски RE и регистр адреса операнда RF.

Была составлена и отлажена микропрограмма командного цикла ЭВМ, обеспечивающая начальный пуск, выборку команд из ОЗУ, выделение адресной части с помощью маски, дешифрацию кода операции через таблицу начальных адресов ПНА и формирование программного признака нуля PZ на основе аппаратного флага $Z$. На языке машинных кодов разработана и выполнена программа циклического логического сдвига 16-разрядного двоичного числа $X$ на $N$ разрядов влево. Работоспособность микропрограммы и программы подтверждена в имитаторе микропрограммируемого микропроцессора.

\end{document}